\section{Метод стрижки дерева}

Деревья решений демонстрируют выраженную склонность к переобучению, особенно при построении до максимальной глубины. Это можно сравнить со студентом, который вместо понимания общего принципа зазубривает конкретные примеры — на знакомых задачах он работает идеально, но при малейшем изменении условий теряется. Это хорошо для тренировки, но плохо для новых данных.

\textit{Стрижка дерева} похожа на то, как опытный садовник сначала дает дереву вырасти, а затем аккуратно подрезает лишние ветви. Сначала строится переобученное дерево, а затем удаляются неоптимальные поддеревья. Этот метод позволяет исследовать более широкое пространство моделей и полностью контроллировать качество обучения модели.

\subsection{Реализация стрижки}

Далее объяснение будет приведено <<на пальцах>>, т.к. непосредственно алгоритмическая реализация несёт в себе добавление новых критериев и потерь, а Вы нам нужны ещё живыми.

\begin{enumerate}
  \item \textbf{Выращиваем максимальное дерево}
  
  Сначала мы позволяем дереву расти полностью, без ограничений. Каждый листовой узел содержит по несколько (или даже по одному) объектов. Это переобученное дерево — наша начальная точка.
  
  \item \textbf{Ищем кандидатов на удаление}
  
  Для каждого узла мы задаёмся вопросом: «А нужна ли эта ветка на самом деле?» Проверяем, что произойдёт, если мы превратим этот узел обратно в лист (то есть удалим всё его поддерево). Как сильно ухудшится точность?
  
  \item \textbf{Удаляем самые бесполезные ветки}
  
  Из всех кандидатов мы выбираем ту ветку, которая даёт минимальный выигрыш для модели. Если её удаление почти не ухудшает результат, но сильно упрощает дерево — убираем её.
  
  \item \textbf{Повторяем процесс}
  
  Продолжаем шаги 2 и 3, пока не получим серию упрощённых деревьев — от очень сложного к очень простому.
\end{enumerate}